\subsection{Denotational Semantics of \GRSync}
\label{sec:semantics:grsync:denot}

\GRSync's semantics is based on prior work on dataflow and synchronous
languages~\cite{DBLP:conf/ifip/Kahn74,DBLP:journals/tecs/BourkePP23}, associating to each identifier
an infinite sequence of values, called \emph{stream}, representing the value the identifier will carry
at every time step. Time steps are represented by a clock ($ck$), a stream of boolean that indicates
the rhythm of other streams and therefore of their computations.

\begin{table}[!ht]
    \centering
    \begin{tabularx}{\textwidth}{|l|X|}\hline
        $v, \, \some{v}, \, \none \in \Val$                           & A value of any type; \ttf{some} and \none represent optional values.                                           \\\hline
        $v_\bot \in \ValBot ::= v \, \vert \, \bot$                   & A value that may be present ($v$) or absent ($\bot$).                                                          \\\hline
        $sty ::= ty \, \vert \, \mayB{ty}$                            & A type, where \kwf{?} indicates an event type.                                                                 \\\hline
        $\tyCtx[x] = sty$                                             & A typing context \tyCtx mapping flow identifiers to their types.                                               \\\hline
        $\G[f]$                                                       & A program context \G mapping component and service names to their definitions.                                 \\\hline
        $s \in \Stream{\ValBot} ::= v_\bot \step s'$                  & A coinductive stream of bottom values.                                                                         \\\hline
        $ck ::= (\true \, \vert \, \false \, \vert \, \bot) \step ck$ & A clock is a stream of booleans.                                                                               \\\hline
        $[s^+]$                                                       & An implicitly finite list of streams $s^+$.                                                                    \\\hline
        $op(s^+)$                                                     & An operation $op$ on values that is implicitly lifted to the streams $s^+$.                                    \\\hline
        % \const{c}{ck}                                                   & A stream of constant $c$ on clock $ck$.                                                                                \\\hline
        % \buffer{v}{s}                                                   & A buffer of stream $s$ initialized by value $v$.                                                                       \\\hline
        % \whenB{s}{ck}                                                   & Filtered version of stream $s$ when values of $ck$ equal $b$.                                                          \\\hline
        % \whenB{H}{ck}                                                   & Filtered version of streams in $H$ when values of $ck$ equal $b$.                                                      \\\hline
        $H[x], \, H[\last{x}]$                                        & A stream context $H$ that maps identifiers to their stream.                                                    \\\hline
        $s_1 \equiv s_2$                                              & Equivalence of the streams $s_1$ and $s_2$.                                                                    \\\hline
        \compDenot{ck}{f}{s_i^+}{s_o^+}                               & The component $f$ is a function from the input streams $s_i^+$ to the output streams $s_o^+$, defined on $ck$. \\\hline
        \eqDenot{H}{ck}{eq^+}                                         & The constraints of equations $eq^+$ are satisfied by the context $H$ on $ck$.                                  \\\hline
        \exprDenot{H}{ck}{e}{s^+}                                     & The expression $e$ is associated with the streams $s^+$ on $ck$.                                               \\\hline
        \patDenot{H}{ck}{\mathfoot{\guarded{pat^+ = s^+}{e_g}}}{ck'}  & The guarded patterns \guarded{pat^+}{e_g} matching the streams $s^+$ are associated with the clock $ck'$.      \\\hline
        \epatDenot{H}{ck}{\guarded{epat^+}{e_g}}{ck'}                 & The guarded event patterns \guarded{epat^+}{e_g} are associated with the clock $ck'$.                          \\\hline
    \end{tabularx}
    \caption[Notations for the denotational semantics of \GRSync.]%
    {Notations for the denotational semantics of \GRSync.}
    \label{tab:semantics:grsync:denot:notations}
\end{table}

We now introduce the notations used in this section to define the denotational semantics. These
notations, which build upon those introduced for the operational semantics in
\Cref{sec:semantics:grsync:opsem}, are summarized in \Cref{tab:semantics:grsync:denot:notations}.

A stream $s$ is an infinite sequence of bottom values, coinductively constructed as
$v \step s'$ (\resp $\bot \step s'$) meaning the first value of $s$ is defined (\resp undefined) and
its next values are those of $s'$. There is no base construction for stream, as they are infinite.
The set of all possible streams of bottom values is denoted by \Stream{\ValBot}.
%
Clocks are streams of booleans, denoted by $ck$, that represent the logical instants of computation
and its rate. We say that a clock \emph{ticks} when its value is \true. More information on clocks
and streams are provided in \Cref{sec:semantics:streams:def,sec:semantics:streams:clock}.
%
Finite lists of streams are denoted by $[s^+]$.
%
Reasoning on streams is done using the coinduction principle~\cite{DBLP:journals/mscs/KozenS17}.
This principle, further explained in \Cref{sec:semantics:streams:coinduction}, provides a framework
for reasoning about infinite structures like streams. Specifically, it allows us to apply the
\emph{coinduction hypothesis} only to streams that have \emph{progressed} and remain \emph{opaque}.
%
Operations on values are implicitly lifted to streams ($op(s^+)$), using the definition
in \Cref{eq:semantics:grsync:denot:aux:lift}. This definition is proved to be well-defined using the
coinduction principle: the construction of $op((v_\bot \step s)^+)$ calls the definition of
$op(s^+)$ (the diminution from $(v_\bot \step s)^+$ to $s^+$ shows progress), where no assumptions
are made about $s^+$ ($s^+$ are opaque), we invoke the coinduction hypothesis to establish that
$op(s^+)$ is well-defined, we then prove that $op((v_\bot \step s)^+)$ is well-defined.
\begin{align}
    op((v \step s)^+)    & = op(v^+) \step op(s^+) &
    op((\bot \step s)^+) & = \bot \step op(s^+)
    \label{eq:semantics:grsync:denot:aux:lift}
\end{align}
%
A context $H$ stores the streams associated with identifiers. For an identifier $x$, $H[x] = s$
indicates that $x$ is associated with the stream $s$.
%
The equivalence of streams $s_1$ and $s_2$ is written $s_1 \equiv s_2$ and states their values are
identical, as defined coinductively by the rule \nameref{eq:semantics:streams:equiv}
(in \Cref{sec:semantics:streams:def}).

Using those notations, we refine the objective of the denotational semantics of \GRSync. It
defines how to construct the streams corresponding to each identifier and expression. It models
equations as constraints on the stream context $H$, and represents causal components as
stream functions.
%
The denotational view abstracts away the control flow of evaluation and focuses on the functional
relationship between input and output streams.
%
The following paragraphs explain in details the mutually recursive rules that define this semantics.

\paragraph{Components}

In the denotational semantics, a causal component $f$ is treated as a function from input streams
$s_i^+$ to output streams $s_o^+$, with both streams sharing a clock $ck$. The function is denoted
by:
\begin{equation*}
    \compDenot{ck}{f}{s_i^+}{s_o^+}.
\end{equation*}

Following rule \nameref{rule:semantics:grsync:denot:comp}, the function maps $s_i^+$ to $s_o^+$ if
there exists a local stream context $H$ associating inputs, outputs and memories to their respective
streams and satisfying the constraints of the equations $eq^+$ of the component $f$.

The determinism of the semantics under causality, stated in
\Cref{thm:semantics:grsync:theorems:determinism}
(in \Cref{sec:semantics:grsync:theorems:determinism}), guarantees that this context $H$ is unique
for the identifiers of $f$. The uniqueness of $H$ ensures that $f$ is a well-defined function.

\begin{center}
    \begin{minipage}{0.9\textwidth}
        \begin{mathpar}
            \inferDENOT
            {
                G[f] = \component{f}{x_i^+}{x_o^+}{eq^+}{x_m^+ = c^+} \\\\
                \left( H[x_i] \equiv s_i \right)^+ \\
                \left( H[x_o] \equiv s_o \right)^+ \\
                \left( H[\last{x_m}] \equiv \buffer{c}{H[x_m]} \right)^+ \\\\
                \eqDenot{H}{ck}{eq^+}
            }
            { \compDenot{ck}{f}{s_i^+}{s_o^+} }
            {Comp}
            \label{rule:semantics:grsync:denot:comp}
        \end{mathpar}
    \end{minipage}
\end{center}

The \buffer{c}{s} function, coinductively defined in \Cref{eq:semantics:grsync:denot:aux:buffer},
creates a stream whose first value is $c$ and the followings are the ones from $s$, shifted so that
it is on the same clock as $s$.
\begin{align}
    \buffer{v_1}{\bot \step s} & = \bot \step \buffer{v_1}{s} &
    \buffer{v_1}{v_2 \step s}  & = v_1 \step \buffer{v_2}{s}
    \label{eq:semantics:grsync:denot:aux:buffer}
\end{align}

Let us illustrate the component semantics on the example \Cref{fig:semantics:example:aeb}.
The denotational semantics of the component \idf{braking\_strat} corresponds to the stream function
denoted by:
\begin{equation*}
    \idf{ck} \vdashSync \idf{braking\_strat}[\idf{p}, \idf{t\_p}, \idf{s}] \Downarrow [\idf{b}]
\end{equation*}
where \idf{ck} is the clock on which the streams \idf{p}, \idf{t\_p}, \idf{s}, and \idf{b} are
defined. Events \idf{p} and \idf{t\_p} are streams of optional values, with occurrences wrapped in
\ttf{some} and absences replaced by \none. Each stream is defined at discrete instants when the
clock ticks, and the output stream \idf{b} results from evaluating the component equations at
every clock tick.
%
For instance, based on the execution from \Cref{fig:semantics:example:aeb:chronogram} with the clock
\idf{ck} constantly \true, the first values of these streams for the application in
\Cref{ex:line:strat} are:%
\setlength{\arraycolsep}{2pt}%
\begin{equation*}
    \begin{array}{rrlllllllllll}
        \idf{p}    & = & \none & \step & \some{20.} & \step & \none & \step & \none & \step & \none     & \cdots \\
        \idf{t\_p} & = & \none & \step & \none      & \step & \none & \step & \none & \step & \some{()} & \cdots \\
        \idf{s}    & = & 32.   & \step & 32.        & \step & 30.   & \step & 25.   & \step & 25.       & \cdots \\
        \idf{b}    & = & \No   & \step & \Soft      & \step & \Soft & \step & \Soft & \step & \No       & \cdots
    \end{array}
\end{equation*}%
\setlength{\arraycolsep}{5pt}%

\paragraph{Equations}

The denotational semantics models equations ($eq^+$) as constraints for the stream context $H$ on
the clock $ck$. The satisfaction of those constraints is denoted by:
$$\eqDenot{H}{ck}{eq^+}.$$

The rule \nameref{rule:semantics:grsync:denot:eqs} states that a stream context $H$
satisfies the constraints of a block of equations ($eq^+$) if it satisfies the constraints on
each individual equations ($eq$).

\begin{center}
    \begin{minipage}{0.9\textwidth}
        \begin{mathpar}
            \inferDENOTEq
            { \left(\eqDenot{H}{ck}{eq}\right)^+ }
            { \eqDenot{H}{ck}{eq^+} }
            {Blck}
            \label{rule:semantics:grsync:denot:eqs}
        \end{mathpar}
    \end{minipage}
\end{center}

For identifiers declarations ($x^+ = e$), the rule \nameref{rule:semantics:grsync:denot:eq:decl}
states that the streams associated with the identifiers $x^+$ in the context $H$ are equivalent
to those associated with the expression $e$.

\begin{center}
    \begin{minipage}{0.9\textwidth}
        \begin{mathpar}
            \inferDENOTEq
            {
                \exprDenot{H}{ck}{e}{s_e^+} \\
                \left(H[x] \equiv s_e\right)^+
            }
            { \eqDenot{H}{ck}{x^+ = e} }
            {Decl}
            \label{rule:semantics:grsync:denot:eq:decl}
        \end{mathpar}
    \end{minipage}
\end{center}

The rule \nameref{rule:semantics:grsync:denot:eq:match} states that a stream context $H$ satisfies
the constraints of a \kwf{match} equation (\match{e}{x^+}{(\branch{pat_k^+}{e_k}{eq_k^+})^+}) if it
satisfies the constraints of each branch-equations $eq_k^+$ on the clock $ck'_k$. This clock
corresponds to associated clock of the guarded patterns \guarded{pat_k^+ = s^+}{e_k}, filtered
when the clocks of the prioritized patterns are \true.%
\footnote{There can be multiple patterns that hold at the same time, like in usual pattern matching,
    only the first successful pattern is taken.}
%
This corresponds to clock filtering and merging in \Lustre. The function \whenB{s}{ck} (\resp
\whenB{H}{ck}), coinductively defined in \Crefrange{eq:semantics:grsync:denot:aux:when:stream1}
{eq:semantics:grsync:denot:aux:when:stream2} (\resp in \Cref{eq:semantics:grsync:denot:aux:when:context}), is
a \Lustre expression used to filter $s$ (\resp the streams in $H$) only when the value of clock $ck$
is equal to $b$.
%
Pattern exhaustiveness ensures the total merge is on the global clock $ck$.
\begin{align}
    \whenB{\bot \step s}{\bot \step ck} & = \bot \step \whenB{s}{ck}
    \label{eq:semantics:grsync:denot:aux:when:stream1}                       \\
    \whenB{v \step s}{b' \step ck}      & = (
    \ttf{if}   \, b' = b \,
    \ttf{then} \, v      \,
    \ttf{else} \, \bot
    ) \step \whenB{s}{ck} \label{eq:semantics:grsync:denot:aux:when:stream2} \\
    \whenB{H}{ck}[x]                    & = \whenB{H[x]}{ck}
    \label{eq:semantics:grsync:denot:aux:when:context}
\end{align}

\begin{center}
    \begin{minipage}{0.9\textwidth}
        \begin{mathpar}
            \inferDENOTEq
            {
                \exprDenot{H}{ck}{e}{s^+} \\
                \left(\patDenot{H}{ck}{\guarded{pat_k^+ = s^+}{e_k}}{ck_k}\right)^+ \\
                ck \equiv \kwf{\Vert}(ck_k)^+ \\
                \left(ck'_k = ck_k \kwf{\&\&} (\kwf{\&\&}(\kwf{!}ck_i)^{i < k})\right)^+ \\
                \left(H^\T_k = \whenT{H}{ck'_k} \right)^+\\
                \left(\eqDenot{H^\T_k}{ck'_k}{eq_k^+}\right)^+
            }
            { \eqDenot{H}{ck}{\match{e}{x^+}{(\branch{pat_k^+}{e_k}{eq_k^+})^+}} }
            {Mtch}
            \label{rule:semantics:grsync:denot:eq:match}
        \end{mathpar}
    \end{minipage}
\end{center}

Similar to the \kwf{match} equations, the rule \nameref{rule:semantics:grsync:denot:eq:when} states
that a stream context $H$ satisfies the constraints of a \kwf{when} equation
\begin{equation*}
    \when{x^+}{x_m^+ = c^+}{(\branch{epat_k^+}{e_k}{eq_k^+})^+}
\end{equation*}
if it satisfies the constraints of each
branch-equations $eq_k^+$ on the clock $ck'_k$, the filtered version of the clock associated with the
guarded event patterns \guarded{epat_k^+}{e_k}, with default values for identifiers not present in equations $eq_k^+$ or in case there is no
event detected: previous values for signals ($x_m^+$ identifiers) and \none for events.
This is why signals has to be initialized in \init{x_m^+ = c^+}.
%
The rule for the \kwf{when} equation relies on the function \const{c}{ck}, coinductively defined in
\Crefrange{eq:semantics:grsync:denot:aux:const1}{eq:semantics:grsync:denot:aux:const3}, that builds
a stream of constant values equal to $c$ on the clock $ck$.
\begin{align}
    \const{c}{\bot \step ck}   & = \bot \step \const{c}{ck} \label{eq:semantics:grsync:denot:aux:const1} \\
    \const{c}{\true \step ck}  & = c \step \const{c}{ck}                                                 \\
    \const{c}{\false \step ck} & = \bot \step \const{c}{ck} \label{eq:semantics:grsync:denot:aux:const3}
\end{align}

\begin{center}
    \begin{minipage}{0.9\textwidth}
        \begin{mathpar}
            \inferDENOTEq
            {
                \left(\epatDenot{H}{ck}{\guarded{epat_k^+}{e_k}}{ck_k}\right)^+ \\\\
                \left(ck'_k = ck_k \kwf{\&\&} (\kwf{\&\&}(\kwf{!}ck_i)^{i < k})\right)^+ \\\\
                \left(H^\T_k = \whenT{H}{ck'_k} \right)^+\\
                \left(\eqDenot{H^\T_k}{ck'_k}{eq_k^+}\right)^+ \\\\
                \left(\forall x_s \in \{x_m^+\} \backslash \defs{eq_k^+}, H^\T_k[x_s] \equiv H^\T_k[\last{x_s}]\right)^+ \\\\
                \left(\forall x_e \in \{x^+\} \backslash (\{x_m^+\} \cup \defs{eq_k^+}), H^\T_k[x_e] \equiv \const{\none}{ck'_k}\right)^+ \\\\
                ck^\F = \whenF{ck}{\kwf{\Vert}(ck_k^+)} \\
                H^\F = \whenF{H}{\kwf{\Vert}(ck_k^+)} \\\\
                \forall x_s \in \{x_m^+\}, H^\F[x_s] \equiv H^\F[\last{x_s}] \\
                \forall x_e \in \{x^+\} \backslash \{x_m^+\}, H^\F[x_e] \equiv \const{\none}{ck^\F} \\
                \left(H[\last{x_m}] \equiv \buffer{c}{H[x_m]} \right)^+
            }
            { \eqDenot{H}{ck}{\when{x^+}{x_m^+ = c^+}{(\branch{epat_k^+}{e_k}{eq_k^+})^+}} }
            {When}
            \label{rule:semantics:grsync:denot:eq:when}
        \end{mathpar}
    \end{minipage}
\end{center}

\paragraph{Expressions}

A \GRSync expression $e$ is associated with the streams of its values $s^+$.
With a stream context $H$, that associates streams to identifiers, on a clock $ck$, this association
is denoted by:
$$\exprDenot{H}{ck}{e}{s^+}.$$

A constant expression $c$ follows the rule \nameref{rule:semantics:grsync:denot:e:const} and is
associated with the stream on clock $ck$ constantly equal to $c$. It relies on the function
\const{c}{ck}, defined in \Crefrange{eq:semantics:grsync:denot:aux:const1}{eq:semantics:%
    grsync:denot:aux:const3}, that builds the stream.

\begin{center}
    \begin{minipage}{0.9\textwidth}
        \begin{mathpar}
            \inferDENOTE
            { }
            { \exprDenot{H}{ck}{c}{\const{c}{ck}} }
            {Cst}
            \label{rule:semantics:grsync:denot:e:const}
        \end{mathpar}
    \end{minipage}
\end{center}

Identifiers ($x$) and memories (\last{x}) are associated with the streams stored in the context
($H$) at their respective names ($H[x]$ and $H[\last{x}]$), as defined in rules \nameref{rule:%
    semantics:grsync:denot:e:ident} and \nameref{rule:semantics:grsync:denot:e:last}.

\begin{center}
    \begin{minipage}{0.9\textwidth}
        \begin{mathpar}
            \inferDENOTE
            { }
            { \exprDenot{H}{ck}{x}{H[x]} }
            {ID}
            \label{rule:semantics:grsync:denot:e:ident}

            \inferDENOTE
            { }
            { \exprDenot{H}{ck}{\last{x}}{H[\last{x}]} }
            {Last}
            \label{rule:semantics:grsync:denot:e:last}
        \end{mathpar}
    \end{minipage}
\end{center}

An event emission (\emit{e}) is associated with the stream corresponding to the inner expression ($e$)
wrapped by the option operator \ttf{some} lifted to streams, as defined by the rule \nameref{rule:%
    semantics:grsync:denot:e:emit}.

\begin{center}
    \begin{minipage}{0.9\textwidth}
        \begin{mathpar}
            \inferDENOTE
            { \exprDenot{H}{ck}{e}{s} }
            { \exprDenot{H}{ck}{\emit{e}}{\some{s}} }
            {Emit}
            \label{rule:semantics:grsync:denot:e:emit}
        \end{mathpar}
    \end{minipage}
\end{center}

Operations ($op(e^+)$) are lifted to operate on streams ($op(s^+)$) as defined by rule
\nameref{rule:semantics:grsync:denot:e:op}. Component applications ($f(e^+)$) are associated with
the output streams of the stream function representing the called component ($f$), based on the
input streams corresponding to the input expressions ($e^+$), as specified by rule
\nameref{rule:semantics:grsync:denot:e:app}.

\begin{center}
    \begin{minipage}{0.9\textwidth}
        \begin{mathpar}
            \inferDENOTE
            { \left(\exprDenot{H}{ck}{e}{s}\right)^+ }
            { \exprDenot{H}{ck}{op(e^+)}{op(s^+)} }
            {Op}
            \label{rule:semantics:grsync:denot:e:op}

            \inferDENOTE
            {
                \left(\exprDenot{H}{ck}{e}{s_i}\right)^+ \\\\
                \compDenot{ck}{f}{s_i^+}{s_o^+}
            }
            { \exprDenot{H}{ck}{f(e^+)}{s_o^+} }
            {App}
            \label{rule:semantics:grsync:denot:e:app}
        \end{mathpar}
    \end{minipage}
\end{center}

\paragraph{Patterns}

Guarded patterns (\guarded{pat^+ = s^+}{e_g}), following the rule
\nameref{rule:semantics:grsync:denot:pats}, are associated with a clock ($ck'$) representing match
success for all pattern ($pat = s$), and the success of the guard~($e_g$). This association is
denoted by:
\begin{equation*}
    \patDenot{H}{ck}{\guarded{pat^+ = s^+}{e_g}}{ck'}.
\end{equation*}

Rule \nameref{rule:semantics:grsync:denot:pats} defines the clock $ck_g$ for the guard's success in
two parts. First, it corresponds to the stream associated with the guard $e_g$, conditioned on the
clock $ck_p$ that represents the success of all pattern matches. Second, it corresponds to a stream
of repeated \false values when the patterns do not match. The overall clock $ck_g$ is then a merge
of these two definitions.

\begin{center}
    \begin{minipage}{0.9\textwidth}
        \begin{mathpar}
            \inferDENOTPat
            {
                \left(\patDenot{H}{ck}{pat = s}{ck_1}\right)^+ \\
                ck_p = \kwf{\&\&}(ck_1^+) \\\\
                \exprDenot{\whenT{H}{ck_p}}{\whenT{ck}{ck_p}}{e_g}{\whenT{ck_g}{ck_p}} \\\\
                \whenF{ck_g}{ck_p} = \const{\false}{\whenF{ck}{ck_p}}
            }
            { \patDenot{H}{ck}{\guarded{pat^+ = s^+}{e_g}}{ck_g} }
            {Tple}
            \label{rule:semantics:grsync:denot:pats}
        \end{mathpar}
    \end{minipage}
\end{center}

A constant pattern ($c = s$) follows rule \nameref{rule:semantics:grsync:denot:pat:const} and
matches when the value of the matched stream ($s$) equals the constant ($c$). The notation $(\cdot
    \kwf{==} c)$ indicates a single-argument function that substitutes the $\cdot$ with the value of
the input.

\begin{center}
    \begin{minipage}{0.9\textwidth}
        \begin{mathpar}
            \inferDENOTPat
            { }
            { \patDenot{H}{ck}{c = s}{(\cdot \kwf{==} c)(s)} }
            {Cst}
            \label{rule:semantics:grsync:denot:pat:const}
        \end{mathpar}
    \end{minipage}
\end{center}

An identifier pattern ($x = s$) always holds, so the associated clock is constantly \true on
the base clock ($ck$), and the stream associated with the identifier in the context ($H[x]$) is
equivalent to the matched stream $s$, as stated by the rule \nameref{rule:semantics:grsync:denot:%
    pat:ident}.

\begin{center}
    \begin{minipage}{0.9\textwidth}
        \begin{mathpar}
            \inferDENOTPat
            { H[x] \equiv s }
            { \patDenot{H}{ck}{x = s}{\const{\true}{ck}} }
            {ID}
            \label{rule:semantics:grsync:denot:pat:ident}
        \end{mathpar}
    \end{minipage}
\end{center}

A wild pattern ($\_ = s$) always holds, as depicted by the rule \nameref{rule:semantics:grsync:%
    denot:pat:wild}.

\begin{center}
    \begin{minipage}{0.9\textwidth}
        \begin{mathpar}
            \inferDENOTPat
            { }
            { \patDenot{H}{ck}{\_ = s}{\const{\true}{ck}} }
            {Wild}
            \label{rule:semantics:grsync:denot:pat:wild}
        \end{mathpar}
    \end{minipage}
\end{center}

\paragraph{Event patterns}

Guarded event patterns (\guarded{epat^+}{e_g}), following the rule \nameref{rule:semantics:grsync:%
    denot:epats}, are associated with a clock ($ck'$) representing occurrence of all event pattern
($epat$), and the success of the guard ($e_g$). This association is denoted by:
$$\epatDenot{H}{ck}{\guarded{epat^+}{e_g}}{ck'}.$$

\begin{center}
    \begin{minipage}{0.9\textwidth}
        \begin{mathpar}
            \inferDENOTEPat
            {
                \left(\patDenot{H}{ck}{epat}{ck_1}\right)^+ \\
                ck_p \equiv \kwf{\&\&}(ck_1^+) \\\\
                \exprDenot{\whenT{H}{ck_p}}{\whenT{ck}{ck_p}}{e_g}{\whenT{ck_g}{ck_p}} \\\\
                \whenF{ck_g}{ck_p} \equiv \const{\false}{\whenF{ck}{ck_p}}
            }
            { \epatDenot{H}{ck}{\guarded{epat^+}{e_g}}{ck_g} }
            {Tple}
            \label{rule:semantics:grsync:denot:epats}
        \end{mathpar}
    \end{minipage}
\end{center}

An event pattern (\detect{x}{y}) follows rule \nameref{rule:semantics:grsync:denot:epat:event}: it
gets the clock ($ck'$) that is \true when the stream associated with the event in the context
($H[y]$) has \ttf{some} values, and states that the stream associated with the receiving identifier
in the context ($H[x]$) is equivalent to $H[y]$ when wrapped by \ttf{some} and filtered on $ck'$.

\begin{center}
    \begin{minipage}{0.9\textwidth}
        \begin{mathpar}
            \inferDENOTEPat
            {
                ck' \equiv \isSome{H[y]} \\
                \whenT{H[y]}{ck'} \equiv \whenT{\some{H[x]}}{ck'}
            }
            { \epatDenot{H}{ck}{\detect{x}{y}}{ck'} }
            {Dtct}
            \label{rule:semantics:grsync:denot:epat:event}
        \end{mathpar}
    \end{minipage}
\end{center}

A rising edge ($e$) is associated with a clock that is \true when the value of the stream
associated with the boolean expression $e$ was \false in the previous time step and \true in
the current one, as defined by rule \nameref{rule:semantics:grsync:denot:epat:edge}.
It relies on the \buffer{c}{s} function, defined in \Cref{eq:semantics:grsync:denot:aux:%
    buffer}, that creates a stream whose first value is $c$ and the followings are the ones from
$s$, shifted so that it is on the same clock as $s$.

\begin{center}
    \begin{minipage}{0.9\textwidth}
        \begin{mathpar}
            \inferDENOTEPat
            {
                \exprDenot{H}{ck}{e}{ck_1} \\
                ck_2 \equiv ck_1 \kwf{\&\&} \kwf{!} \buffer{\false}{ck_1}
            }
            { \epatDenot{H}{ck}{e}{ck_2} }
            {Edge}
            \label{rule:semantics:grsync:denot:epat:edge}
        \end{mathpar}
    \end{minipage}
\end{center}

\paragraph{Summary}
This section detailed the denotational semantics for \GRSync, grounding it in the classical theory
of synchronous languages. In this model, identifiers are interpreted as infinite streams of values,
whose temporal evolution is governed by clocks. Components are defined as functions that transform
input streams into output streams, while equations impose constraints on these streams. We
demonstrated how constructs like \kwf{match} and \kwf{when} are handled through a clock calculus,
using clock filtering and merging. This stream-based formalism provides a high-level, declarative
view of program semantics, abstracting away from operational details to focus on the functional
relationships between dataflows.
